\documentclass[12pt,letterpaper]{article}
\usepackage[utf8]{inputenc}
\usepackage[spanish]{babel}
\usepackage[margin=2.5cm]{geometry}
\usepackage{xcolor}
\usepackage{hyperref}

\begin{document}

\begin{center}
    \Large \textbf{Curso de Seguridad Informática} \\
    \large \textbf{Suite de Auditoría en Python} \\
    \vspace{0.3cm}
    \normalsize Prof. César Rodríguez \\
\end{center}

\vspace{0.5cm}

\noindent \textbf{Fecha:} \today \\
\textbf{Asunto:} Tareas y plazos para la Fase III (Aplicaciones Web y Reportes) \\
\textbf{Plazo de entrega:} Semanas 8 a 10

\vspace{0.5cm}

\noindent Estimados estudiantes,

\vspace{0.3cm}

Espero que se encuentren muy bien. Tras los excelentes resultados obtenidos en la fase anterior, me dirijo a ustedes para dar inicio formal a la \textbf{Fase III: Aplicaciones Web y Reportes}, la etapa final de nuestro proyecto integrador. 

Durante las \textbf{semanas 8, 9 y 10} de nuestro calendario, nos enfocaremos en el análisis de vulnerabilidades en aplicaciones web y la consolidación de nuestra suite. La asignación de tareas por grupo será la siguiente:

\begin{itemize}
    \item \textbf{Grupo 1:} Web Crawler, Mapeo y Fuerza Bruta de Directorios (\texttt{web\_crawler.py})
    \item \textbf{Grupo 2:} Módulo Principal (Orquestador) y Generación de Reportes (\texttt{auditoria.py})
    \item \textbf{Grupo 3:} Escáner de Vulnerabilidades SQLi (\texttt{vuln\_sqli.py})
    \item \textbf{Grupo 4:} Escáner de Vulnerabilidades XSS y LFI/RFI (\texttt{vuln\_xss\_lfi.py})
\end{itemize}

\vspace{0.3cm}
\noindent \textbf{Instrucciones y Entregables:}
\begin{enumerate}
    \item \textbf{Material Teórico:} En esta etapa nos basaremos en los Capítulos 11 y 12 del libro guía. Es fundamental comprender la estructura de las peticiones HTTP y la lógica de las vulnerabilidades antes de codificar.
    \item \textbf{Contrato de Datos:} Como siempre, los resultados de los escáneres de los grupos 1, 3 y 4 deben ajustarse obligatoriamente al esquema \texttt{schema\_resultados.json}. Esto es vital para que el Grupo 2 pueda unificar la data y generar los reportes HTML/CSV sin errores.
    \item \textbf{Entrega (Pull Request):} El código finalizado deberá ser enviado mediante un \textbf{Pull Request} en GitHub antes de finalizar la \textbf{Semana 10}.
\end{enumerate}

\vspace{0.3cm}
\noindent \textbf{Recordatorios Críticos de Integración:}
\begin{itemize}
    \item \textbf{Carpetas:} Todo el desarrollo de módulos (Grupos 1, 3 y 4) debe guardarse en la carpeta \texttt{modulos/}.
    \item \textbf{Exclusividad del Grupo 2:} El Grupo 2 es el \textbf{único} equipo autorizado para modificar el orquestador principal \texttt{auditoria.py} durante esta etapa.
\end{itemize}

¡Mucho éxito en esta recta final! Estoy ansioso por ver nuestra Suite de Auditoría completamente funcional e integrada.

\vspace{1cm}
\noindent Atentamente, \\
\vspace{0.5cm} \\
\textbf{Prof. César Rodríguez}

\end{document}